\documentclass[12pt]{article}

% Spacing between paragrphes
\usepackage[skip=10pt plus1pt, indent=15pt]{parskip}

% Language setting
% Replace `english' with e.g. `spanish' to change the document language
\usepackage[T1]{fontenc}
\usepackage[french]{babel}

% For figures and tables
\usepackage{graphicx}
\graphicspath{ {img/} }
\usepackage{array}

% Glossary
\usepackage[toc, acronym]{glossaries}
\setacronymstyle{long-short-desc}
\makeglossaries
% ==========================
% ======== Acronyms ========
% ==========================

\newacronym{vn}{VN}{Visual Novel}
\newacronym{ia}{I.A.}{Intelligence Artificielle}
\newacronym{json}{JSON}{JavaScript Object Notation}



% ============================
% ======== Glossaries ========
% ============================

\newglossaryentry{opensource}{
    name={open source},
    description={Modèle de développement dans lequel le code source est public et peut être consulté, modifié et redistribué par tous, selon les conditions de sa licence.}
}

% Set page size and margins
% Replace `letterpaper' with `a4paper' for UK/EU standard size
\usepackage[a4paper,top=2cm,bottom=2cm,left=3cm,right=3cm,marginparwidth=1.75cm]{geometry}

% Useful packages
\usepackage{amsmath}
\usepackage{graphicx}
\usepackage{booktabs}
\usepackage[]{mdframed}
\usepackage{datetime2}
\usepackage{titlesec}
\usepackage{subcaption}
\usepackage{minted}
\usepackage{tablefootnote}


\usepackage{tikz}
\usetikzlibrary{arrows.meta, positioning, fit, backgrounds}
\usepackage{xcolor}


% ----------------------------------------------------------

% Couleurs
\definecolor{clrInput}   {HTML}{D6EAF8}
\definecolor{clrProcess} {HTML}{FEF9E7}
\definecolor{clrToken}   {HTML}{D5F5E3}
\definecolor{clrOutput}  {HTML}{FADBD8}
\definecolor{clrBorder}  {HTML}{2C3E50}
\definecolor{clrArrow}   {HTML}{1A252F}
\definecolor{clrNote}    {HTML}{7F8C8D}

\definecolor{clrAuteur} {HTML}{D6EAF8}
\definecolor{clrBuild}  {HTML}{FEF9E7}
\definecolor{clrCloud}  {HTML}{D5F5E3}
\definecolor{clrPlayer} {HTML}{FADBD8}
\definecolor{clrAlt}    {HTML}{EBE2F5}
\definecolor{clrBorder} {HTML}{2C3E50}
\definecolor{clrArrow}  {HTML}{1A252F}
\definecolor{clrNote}   {HTML}{7F8C8D}

% ----------------------------------------------------------


\usepackage{pgfplots}
\pgfplotsset{compat=1.18}


\usepackage[style=ext-authoryear, backend=biber]{biblatex}
\addbibresource{references.bib} % note the .bib is required

\usepackage[colorlinks=true, allcolors=blue]{hyperref}


% ----------------------------------------------------------

\newcommand{\ns}{NamelessStory}
\newcommand{\nsit}{\textit{\ns}}
\newcommand{\nsits}{\nsit\space}

% ----------------------------------------------------------


% ----------------------------------------------------------

\renewcommand{\listingscaption}{Extrait de code}
\renewcommand{\listoflistingscaption}{Liste des extraits de code}

% ----------------------------------------------------------


% ----------------------------------------------------------

\titleclass{\subsubsubsection}{straight}[\subsection]

\newcounter{subsubsubsection}[subsubsection]
\renewcommand\thesubsubsubsection{\thesubsubsection.\arabic{subsubsubsection}}

\titleformat{\subsubsubsection}
  {\normalfont\normalsize\bfseries}{\thesubsubsubsection}{1em}{}
\titlespacing*{\subsubsubsection}
{0pt}{3.25ex plus 1ex minus .2ex}{1.5ex plus .2ex}
\makeatletter
\def\toclevel@subsubsubsection{4}
\def\l@subsubsubsection{\@dottedtocline{4}{7em}{4em}}
\makeatother
\setcounter{secnumdepth}{4}
\setcounter{tocdepth}{4}

% ----------------------------------------------------------


\DTMsetdatestyle{ddmmyyyy}

\title{Mémoire de Bachelor - Da Silva Pinto Kevin}
\author{Da Silva Pinto Kevin}

\begin{document}

\begin{titlepage}
    \begin{center}
        \vspace*{1cm}

        \textbf{Da Silva Pinto Kevin}

        \#12-13\#\#\#\#

        \vspace*{1cm}
            
        BACHELOR OF SCIENCE

        WEB
            
        \vspace{5cm}

        \begin{mdframed}
            \begin{center}
                \textbf{Dans quelle mesure une architecture de moteur de visual novel web basée sur un format de script déclaratif permet-elle à des auteurs non-développeurs de créer du contenu interactif?}
            \end{center}
        \end{mdframed}
            
        \vspace{1.5cm}

        6FSC1WEBB102 - Travail de bachelor

        Référents: Jenna Luthi, Julien Ramirez
        
        \vfill

        SAE Genève

        \today
            
    \end{center}
\end{titlepage}

\addcontentsline{toc}{section}{Préface}
\section*{Préface}

C'est en me baladant sur Steam que j'ai découvert les \acrlong{vn}s, en tombant par hasard sur \textit{Higurashi When They Cry Hou - Ch.1 Onikakushi}. Ce qui m'a frappé, c'est la puissance narrative que ce genre peut atteindre avec des moyens apparemment simples: des images fixes, du texte, de la musique. Depuis, j'ai continué à explorer ce genre, avec des titres comme \textit{Marco et le dragon-galaxie} ou \textit{Find Love or Die Trying}, et mon envie de créer ma propre histoire n'a fait que grandir.

Lorsque je me suis penché sur les outils existants, j'ai rapidement déchanté. RenPy impose l'apprentissage d'un langage de script qui lui est propre. Twine demande une installation et une prise en main qui rebutent quiconque veut simplement écrire une histoire. Je ne voulais pas apprendre un énième outil, je voulais me concentrer sur le récit.

C'est cette frustration qui est à l'origine de \textit{NamelessStory}. L'idée initiale était simple, créer une \gls{librairie} \gls{reactjs} personnelle pour faire tourner mon propre \acrlong{vn}. Mais au fil du développement, le projet a évolué vers quelque chose de plus ambitieux, un moteur \gls{opensource}, accessible depuis un navigateur, pensé pour que n'importe qui disposant de notions de base en \acrshort{json} puisse créer et publier son histoire sans installer quoi que ce soit.

Ce mémoire documente cette évolution, d'une idée personnelle vers un outil conçu pour être partagé.
\clearpage

\addcontentsline{toc}{section}{Remerciements}
\section*{Remerciements}
\clearpage

\addcontentsline{toc}{section}{Résumé}
\section*{Résumé}
\clearpage

% Table des matières ====
\tableofcontents
\clearpage
% =======================

% Lists =================
\listoffigures
\listoftables
\listoflistings
% =======================

% Glossaire =============
\printglossaries
\clearpage
% =======================



% ==========================
% ======== CHAPTERS ========
% ==========================



\section{Introduction}

\subsection{Introduction}




\subsection{Idée créative}




\subsection{Objectifs}

L'objectif de ce travail est d'étudier la faisabilité d'un moteur de \gls{vn} reposant sur un fichier \gls{json} comme source de contenu, rendu directement dans un navigateur web sans nécessiter d'installation ni de compilation.


\subsection{Positionnement}




\subsection{Groupe cible}



\section{Etat de l'art}

\subsection{Le genre Visual Novel}

Le \gls{vn} est un genre de jeu vidéo interactif qui combine narration, visuels, musiques et parfois doublage vocal. Le genre est historiquement ancré dans la culture japonaise, avec des centaines de titres produits chaque année, et a progressivement gagné en popularité en Occident, notamment grâce à la montée en popularité des adaptations animées dans les années 2000.

Sur le plan technique, un \gls{vn} repose sur plusieurs mécaniques fondamentales: une narration textuelle qui progresse par clics, des images statiques représentant les personnages et les décors, des choix de dialogues, et une emphase marquée sur le design sonore et la bande originale.

Camingue, Carstensdottir et Melcer ont relevé que neuf des trente définitions académiques du genre mentionnent la présence de narration à embranchements. Ces structures peuvent aboutir à des expériences très variées: \textit{Steins;Gate 0} (2015) propose six fins différentes, tandis que \textit{428:Shibuya Scramble} (2008) peut mener à pas moins de 85 conclusions distinctes. À l'inverse, 18\% des titres recensés ne comportent aucun embranchement narratif. \cite{camingue_what_2021}

Du point de vue du marché, la période 2019-2024 a connu une croissance de la popularité du genre, porté par l'essor des plateformes de distributions comme Steam et les boutiques d'applications mobiles. Sur la plateforme Steam, nous comptons 16 985 jeux du genre sur Steam en Suisse\footnote{Nombre récupéré au 31 mars 2026}. Le marché mondial du \gls{vn} était estimé à 139 millions de dollars en 2023 et devrait atteindre 620 millions d'ici 2032, avec un taux de croissance annuel moyen de 8,5\%. \cite{noauthor_visual_nodate}



\subsection{Moteur de Visual Novels existants}

Les moteurs de jeu-videos sont des outils permettant de faciliter la création de ces derniers.

Nous pouvons parler de RenPy, un moteur de \gls{vn} utilisés par des milliers de personnes, offrant une interface graphique facile à prendre en main. RenPy permet le développement à destination des ordinateurs comme des téléphones portables.



\subsection{Approches no-code / low-code}

L'objectif de ce type de programmation est de faciliter l'entrée dans la programmation. Ou de permettre à l'aide d'outils graphiques de s'abstraire de l'écriture de code.

L'approche no-code (traduction: pas de code) et low-code (traduction: peux de code) sont très similaires, mais reste différentes. L'une demande toujours quelques connaissances en code ou du moins cela pourrait être utile, tandis que l'autre ne demande aucune connaissances et peut être utilisé directement par tous.


% ================
% Autres chapitres
% ================

\section{Méthodologie}

\subsection{Vue d'ensemble}

La réalisation de \ns s'inscrit dans une démarche de développement logiciel structurée en cycles itératifs. L'objectif central étant de produire un moteur fonctionnel, accessible à des auteurs non-développeurs, la méthode retenue devrait permettre d'ajuster progressivement le périmètre fonctionnel en fonction des contraintes rencontrées: techniques, de format, et de temps.

Le processus global s'organise autour de trois axes: un cycle de développement itératif basé sur le modèle PDCA, une priorisation des fonctionnalités dont les objectifs seront définis au chapitre \ref{sec:objMediaProject}, et une validation fonctionnelle continue à chaque itération.


\subsection{Cycle de développement itératif (PDCA)} \label{sec:PDCA}

Le modèle PDCA (\textit{Plan}, \textit{Do}, \textit{Check}, \textit{Act}), formalisé par W. Edwards Deming dans le cadre de l'amélioration continue des processus \parencite{deming_out_1986}, constitue le cadre méthodologique de ce projet. Chaque itération de développement suit ce cycle en quatre phases:

\begin{table}[H]
    \centering
    \caption{Définition du cycle PDCA utilisée}
    \renewcommand{\arraystretch}{1.5}
    \begin{tabular}{p{1.5cm} p{11cm}}
        \hline
        \textbf{Nom} & \textbf{Description} \\
        \hline
        \textbf{Plan} & Définition des fonctionnalités à implémenter pour l'itération, en référence aux objectifs établis au chapitre. \\
        \addlinespace
        \textbf{Do} & Implémentation des fonctionnalités planifiées. \\
        \addlinespace
        \textbf{Check} & Évaluation du résultat par vérification fonctionnelle manuelle et, le cas échéant, par validation automatique via le \gls{parseur} \gls{zod}. \\
        \addlinespace
        \textbf{Act} & Ajustement des objectifs de l'itération suivante en fonction des écarts constatés. \\
        \hline
    \end{tabular}
    \label{tab:defPDCA}
\end{table}

Ce modèle a été retenu pour sa capacité à intégrer les retours d'évaluation de manière continue, sans nécessiter de redéfinir l'ensemble du périmètre à chaque étape.

\clearpage

Le développement s'est déroulé en trois étapes d'itérations successives, décrites en détail au chapitre \ref{chap:MediaProject}:

\begin{table}[H]
    \centering
    \caption{Objectif par itération}
    \renewcommand{\arraystretch}{1.5}
    \begin{tabular}{p{0.5cm} p{2.7cm} p{9.5cm}}
        \hline
        \textbf{N°} & \textbf{Objectif} & \textbf{Description} \\
        \hline
        1 & Fonctionnement du moteur & Implémentation des mécaniques fondamentales (affichage du texte, gestion des scènes, gestion des choix) et vérification du comportement attendu dans le navigateur. \\
        \addlinespace
        2 & Optimisation & Amélioration des performances internes, notamment par l'utilisation de \glspl{webworker} pour déléguer certains traitements hors du fil d'exécution principal, et refactorisation du code pour en augmenter la maintenabilité. \\
        \addlinespace
        3 & Sécurité et robustesse & Intégration de la \gls{librairie} \gls{zod} pour valider à l'exécution la conformité des fichiers \gls{json} au \gls{schema} attendu, et pour avertir l'utilisateur en cas de fichier de sauvegarde corrompu ou invalide. \\
        \hline
    \end{tabular}
    \label{tab:objectifs_iterations}
\end{table}


\subsection{Critères d'évaluation et validation fonctionnelle}

En l'absence de tests formels automatisés, la validation de chaque fonctionnalité s'est appuyée sur deux approches complémentaires.


\subsubsection{Vérification fonctionnelle manuelle}

Chaque fonctionnalité implémentée a été testée directement dans le navigateur, en condition réelle, à partir de fichiers \gls{json} représentatifs. Les cas de tests couvrent  à la fois des scénarios valides et des scénarios d'erreur intentionnellement construits (\textit{negative testing}\footnote{\textbf{Negative testing} (aussi appelé \textit{unhappy path testing}): approche de test visant à vérifier le comportement du programme face à des entrées inattendues, invalides ou erronées, par opposition au \textit{positive testing} qui couvre le comportement attendu en conditions normales. Hora définit les tests négatifs comme ceux qui vérifient «l'exécution anormale, comme les cas exceptionnels» \parencite{hora_test_2024}. Dans ce projet, les tests négatifs ont consisté à fournir au moteur des fichier \gls{json} syntaxtiquement invalides, des champs manquants ou des valeurs de mauvais type.}): suppression de virgules ou de guillemets pour casser la \gls{syntaxe} \gls{json}, omission de champs obligatoires attendus par le moteur (comme le nom du personnage parlant), et injection de valeurs de mauvais type (par exemple, un booléen à la place d'une chaîne de caractères). Ces fichiers de tests sont disponibles dans le dépôt GitHub du projet. Le critère de validation est la conformité du comportement observé avec le comportement attendu défini lors de la phase \textit{Plan} de chaque itération.


\subsubsection{Validation par Zod} \label{sec:zod}

La \gls{librairie} \gls{zod}, intégrée lors de la troisième itération, valide à l'exécution la structure du fichier \gls{json} fourni par l'auteur. En cas d'écart par rapport au \gls{schema} défini: champ manquant, type incorrect, structure invalide, \gls{zod} retourne un message d'erreur explicite indiquant précisément la source du problème, ce qui constitue un mécanisme de détection d'erreurs déterministe.


\subsection{Considérations sur la sécurité et la robustesse}

Deux préoccupations ont guidé les décisions de développement au-delà de la simple fonctionnalité.


\subsubsection{Robustesse du code}

L'utilisation de \gls{ts} comme surcouche typée de \gls{js} permet de détecter à la compilation un ensemble d'erreurs de type qui seraient silencieuses en \gls{js} pur. Combiné à la validation \gls{zod} à l'exécution, ce dispositif forme une double barrière contre les états invalides du moteur. \parencite{noauthor_starting_nodate}


\subsubsection{Sécurité du code}

Le moteur opère entièrement côté client, dans le navigateur, sans serveur \gls{backend}. Cette architecture supprime une surface d'attaque importante: il n'y a ni base de données exposée, ni gestion d'authentification, ni transmission de données utilisateur vers un serveur tiers. Le contenu narratif est statique et ne transite pas par un serveur applicatif. La validation des fichiers de sauvegarde par \gls{zod} permet par ailleurs d'empêcher l'injection de données malformées susceptibles de mettre le moteur dans un état indéterminé.


\subsection{Objectifs du Media Project} \label{sec:objMediaProject}

L'objectif principal de ce projet est d'offrir une alternative entièrement textuelle aux moteurs de \gls{vn} existants, facile à utiliser et à mettre en ligne. Contrairement aux solutions actuelles, telles que RenPy ou Twine, qui nécessitent l'apprentissage d'un langage de script dédié ou l'installation d'un environnement de développement, \nsits permet à tout auteur disposant de connaissances de base en \gls{json} de créer et de publier un \gls{vn} fonctionnel.

En proposant une solution \gls{opensource}, le projet laisse la possibilité aux personnes possédant des connaissances en développement de modifier et d'étendre les fonctionnalités selon leurs besoins. L'\gls{opensource} permet également à toute personne intéressée de contribuer activement au développement du moteur.

Les objectifs sont définis selon la méthode de priorisation \gls{moscow} \parencite{clegg_case_1994}:

\begin{table}[H]
    \centering
    \caption{Objectifs du Media Project selon la priorisation MoSCoW}
    \renewcommand{\arraystretch}{1.5}
    \begin{tabular}{p{2cm} p{4.5cm} p{6cm}}
        \hline
        \textbf{Priorité} & \textbf{Fonctionnalité} & \textbf{Description} \\
        \hline
        Must & Effet machine à écrire & Affichage du texte caractère après caractère, simulant une machine à écrire. \\
        \addlinespace
        Must & Gestion des scènes & Navigation entre scènes, permettant les changements de décor et la progression narrative. \\
        \addlinespace
        Must & Gestion des choix & Présentation de choix au joueur, orientant la narration vers des routes différentes. \\
        \addlinespace
        Should & Gestion des saisies & Stockage d'une valeur saisie par le joueur, réutilisable dans le récit via une clé. \\
        \addlinespace
        Can & Fonctionnalités additionnelles & Toute autre fonctionnalité envisageable selon le temps disponible (sons, images, transitions). \\
        \hline
    \end{tabular}
    \label{tab:objectifs_moscow}
\end{table}

\clearpage


\subsection{Limites méthodologiques}

Cette approche présente une limite principale: l'absence de tests automatisés (\textit{unit tests}\footnote{\textbf{Test unitaire} (\textit{unit test}): test automatisé vérifiant le comportement d'une unité de code isolée. (\cite{beck_test-driven_2002};\cite{aniche_effective_2022})}, \textit{integration tests}\footnote{\textbf{Test d'intégration} (\textit{integration test}): test automatisé vérifiant que plusieurs unités ou composants fonctionnent correctement ensemble. \parencite{aniche_effective_2022}}). Le développement a privilégié la mise en fonctionnement progressive des fonctionnalités sur la mise en place d'une suite de tests formelle. Cette absence implique que la couverture des cas limites dépend de l'exhaustivité des tests manuels réalisés et que des régressions introduites lors d'une itération pourraient passer inaperçues. Cette limite sera reprise dans l'évaluation critique au chapitre \ref{chap:resultats}.

\section{Media project}

\subsection{Objectifs du media project}

L'objectif du media, est d'offrir une alternative entièrement textuelle aux moteurs de \gl{vn} facile à utiliser et à mettre en ligne.

En proposant une solution \gls{opensource}, nous laissons la possibilité aux personnes possédant des connaissances en développement de modifier et de rajouter des fonctionnalités à leurs projets. De plus, l'\gls{opensource} nous permet de laisser toutes personnes le souhaitant, de participer au développement du projet.



\subsection{Contraintes et limites}

Les moteurs de jeux, afin de rendre ces derniers jouables, compiles le code et le transforme en une version exécutable. Dans le cadre de ce média, nous allons réalisé un outil utilisant des technologies web, nous permettant de voir notre résultat sur n'importe quel navigateur internet. Nous expliquerons également comment mettre en ligne facilement et gratuitement les réalisations.

Une autre contrainte évoqué plus tôt, est de rendre le texte lisible et compréhensible autant pour les humains que pour les ordinateurs. Pour rendre cela possible, nous utiliserons le format \gls{json}. Ce format d'écriture clé/valeur, permet de stocker du texte attribué à des clés.

Pour des raisons de temps à attribué, pour les fonctionnalités nous nous contenterons de celles principales concernant les \gls{vn} qui sont les suivantes.


\begin{table}[htbp]
    \centering
    \caption{Fonctionnalités de bases du media project}
    \renewcommand{\arraystretch}{1.5}
    \begin{tabular}{p{4cm} p{9cm}}
        \hline
        \textbf{Fonctionnalité} & \textbf{Description} \\
        \hline
        Effet machine à écrire
        &
        Simule une machine à écrire, affichant le texte caractère après caractère.
        \\
        \addlinespace
        Gestion des scènes
        &
        Toutes les histoires sont divisées en scènes, le but ici est de permettre de naviguer d'une scène à une autre, d'effectuer un changement de décor.
        \\
        \addlinespace
        Gestion de saisies du joueur
        &
        Mettre un système permettant de stocké une saisie du joueur avec une valeur clé, facile à réutilisé.
        \\
        \addlinespace
        Gestion des choix
        &
        Permettre aux joueurs de choisir une réponse parmi plusieurs, les faisant changer de routes.
        \\
        \hline
    \end{tabular}
\end{table}

D'autres fonctionnalités pourraient être mises en place en fonction du temps, mais pour ce qui est des tests, nous allons nous focaliser sur ces fonctionnalités citées ci-dessus.

\section{Résultats} \label{chap:resultats}

\subsection{Résultat}

% TODO: write here

\begin{table}[htbp]
    \centering
    \caption{Réalisation des objectifs MoSCoW}
    \renewcommand{\arraystretch}{1.5}
    \begin{tabular}{p{1.7cm} p{4cm} p{7cm}}
        \hline
        \textbf{Priorité} & \textbf{Fonctionnalité} & \textbf{Atteint?} \\
        \hline
        Must & Effet machine à écrire & \textbf{Oui}: \\
        \addlinespace
        Must & Gestion des scènes & \textbf{Oui}: \\
        \addlinespace
        Must & Gestion des choix & \textbf{Oui}: \\
        \addlinespace
        Should & Gestion des saisies & \textbf{Oui}: \\
        \addlinespace
        Can & Fonctionnalités additionnelles & Voici les quelques fonctionnalités additionnelles qui ont été ajoutées:  \\
        \hline
    \end{tabular}
    \label{tab:resultats_moscow}
\end{table}

% TODO: write here


\subsection{Évaluation et auto-réflexion}

\subsubsection{Évaluation du Media Project}

% TODO: write here [forces/faiblesses techniques, réponse nuancée à la problématique]


\subsubsection{Réflexion personnelle}

% TODO: write here [organisation, gestion du temps, ce que tu ferais différemment]

\section{Conclusion} \label{chap:conclusion}

\subsection{Synthèse}

% TODO: write here


\subsection{Réponse à la problématique et bilan critique}

% TODO: write here


\subsection{Pistes de développement du Media Project}

% TODO: write here


\subsection{Pistes de recherche}

% TODO: write here



% ==========================
% ==========================
% ==========================



\clearpage

\addcontentsline{toc}{section}{Références}

\printbibliography

\end{document}